\section{统一视角：Pre-training、Post-training 与 RL 的关系}

讲者用 Lego 比喻解释三者关系，这是本讲中最具教学性的部分之一。Pre-training 像学习积木组件与统计结构；Post-training 像根据任务目标组装形态；RL 像在反馈中不断拆改并优化结构。

\textit{``RL is a learning concept; it does not have to be attached only to post-training.''} 这一表述实际上在打破常见分工壁垒。RL 可以服务 pre-training 目标，也可以与 next-token prediction 形成更连续的学习框架。

\begin{knowledgebox}{Lego 比喻的工程映射}
\begin{itemize}
    \item Pre-training：学习 ``哪些砖块存在，如何统计共现''；
    \item Post-training：按任务规范搭建可用结构；
    \item RL：依据反馈迭代修改结构，提高目标匹配度。
\end{itemize}
\end{knowledgebox}

讲者还指出团队里常存在 ``pre-train team vs post-train team'' 的割裂。随着模型能力提升，这种人为边界可能变得低效，未来可能更接近 ``NTP + RL'' 的统一训练叙事。

\begin{importantbox}{统一训练范式的现实意义}
若预训练与后训练共享更多目标和反馈接口，数据回流、评估标准、基础设施可以复用，整体研发效率和一致性都会提升。
\end{importantbox}

\begin{warningbox}{统一不等于混在一起训练}
没有清晰接口的 ``全混合训练'' 反而会导致调试困难。合理路径是先统一评价协议，再逐步统一优化流程。
\end{warningbox}

\subsection{本章小结}
Pre-training 与 Post-training 的边界正在变薄。RL 更像跨阶段的方法学，而非某个固定训练阶段的专属工具。

